\section{Introduction}\label{sec:intro}

A partition of $n$ into distinct parts is a \emph{non-Rascoe partition} if its
number of parts is not itself one of the parts. Let $b(n)$ be the number of
non-Rascoe partitions of $n$, and let $b_1(n)$ be the number of partitions of
$n$ into distinct parts in which the number of parts plus one is not a part.
The empty partition gives $b(0)=b_1(0)=1$. For example, the partitions of $5$
into distinct parts are $5$, $4+1$ and $3+2$; the first contains its number of
parts, so $b(5)=2$. These functions were introduced by Dixit, Kumar and
Srivastava \cite{DKS}, who related them to Rogers--Ramanujan series and
conjectured that $b(29k+21)\equiv0\pmod4$ for $k\ge1$, $29\nmid k$
\cite[Conjecture~1]{DKS}. We prove this conjecture as a special case of the
following criterion. For a prime $p$ and a positive integer $D$, $v_p(D)$
denotes the exponent of $p$ in $D$.

\begin{theorem}\label{thm:main-valuation}
Let $p$ be a prime with $p\equiv21$ or $29\pmod{40}$, and let $n\ge0$. Then
\[
  v_p(40n+1)\ \text{odd}\ \Longrightarrow\ b(n)\equiv0\pmod4,\qquad
  v_p(40n+9)\ \text{odd}\ \Longrightarrow\ b_1(n)\equiv0\pmod4.
\]
In particular,
\begin{align*}
  b(29k+21)&\equiv0\pmod4&&(k\ge1,\ 29\nmid k),\\
  b_1(29k+15)&\equiv0\pmod4&&(k\ge0,\ k\not\equiv6\pmod{29}).
\end{align*}
\end{theorem}

For the prime $29$ we also obtain exact scaling relations, with no condition
on $n$.

\begin{theorem}\label{thm:main-scaling}
For every integer $n\ge0$,
\begin{align}
  b(841n+21)+b(n)&\equiv0\pmod4,\label{eq:scaling-b}\\
  b_1(841n+189)+b_1(n)&\equiv0\pmod4.\label{eq:scaling-b1}
\end{align}
\end{theorem}

For $n\ge1$, congruence \eqref{eq:scaling-b} is Cuthbertson's Conjecture~5.2
\cite{Cuthbertson}; the case $n=0$ also holds. Congruence \eqref{eq:scaling-b1}
is the companion statement for $b_1$. For every prime in \cref{thm:main-valuation},
\cref{cor:scaling-nonsquare} gives a similar scaling relation, provided the
associated number $40n+1$ or $40n+9$ is not a perfect square.

\subsection*{The cone count}
Both partition functions are controlled modulo $4$ by the following lattice
point count. For a positive integer $D\equiv1$ or $9\pmod{40}$ let
\[
  C(D)=\#\{(x,u,v)\in\Z_{>0}^3:\ 10u<x,\ x^2+40uv=D\},
\]
and put $C(D)=0$ if $D$ is not a positive integer. The equation forces
$x\equiv\pm1\pmod{10}$ if $D\equiv1\pmod{40}$ and $x\equiv\pm3\pmod{10}$ if
$D\equiv9\pmod{40}$.

\begin{theorem}\label{thm:cone-parity}
Let $p\equiv21$ or $29\pmod{40}$ be prime. If $D$ is a positive integer that is
not a perfect square, $D\equiv1$ or $9\pmod{40}$, and $p\mid D$, then
\begin{equation}\label{eq:cone-parity}
  C(D)\equiv C(D/p^2)\pmod2 .
\end{equation}
Consequently $C(D)$ is even whenever $v_p(D)$ is odd.
\end{theorem}

\subsection*{How the proof goes}
\Cref{sec:reduction} reduces the generating functions of $b$ and $b_1$ modulo $4$
to a Hecke--Rogers series, whose coefficients vanish at the relevant indices,
plus twice the generating function of $C$. So everything comes down to the
parity of $C(D)$.

For nonsquare $D$ (\cref{sec:involution,sec:certificate,sec:nonsquare}) we
enlarge the cone to the finite set of positive solutions of $x^2+2yz=D$. This
set carries Jackson's involution, a coordinate swap, and two rescalings, and a
sign built from the Legendre symbol modulo $p$ changes in a controlled way
under each of them. A finite identity, valid at every point of the set,
expresses the indicator of the cone as a sum of terms that cancel in pairs
under these maps. The identity depends only on residues modulo $20$ and $200$
and on finitely many linear inequalities, so it can be checked completely by
computer; the data and the checking program are part of the proof.

For square $D$, Jackson's involution has fixed boundaries and the argument
breaks down. Instead (\cref{sec:divisor,sec:recurrence,sec:scaling}) we write
$C(r^2)$ through a divisor sum, express that sum with Lambert series, and reduce
the problem at $p=29$ to the vanishing modulo $2$ of an explicit modular form of
weight two. Sturm's theorem reduces this to finitely many coefficients, which
are computed exactly.

\Cref{sec:verification} explains how to run the two computer verifications,
whose programs are printed in the appendices.
